\uuid{MYps}
\exo7id{1842}
\titre{exo7 1842}
\auteur{legall}
\organisation{exo7}
\datecreate{2003-10-01}
\isIndication{false}
\isCorrection{false}
\chapitre{Extremum, extremum lié}
\sousChapitre{Extremum, extremum lié}
\cpge{mpsi-mp2i.fonctions-deux-variables.b, pcsi.fonctions-deux-variables.b, ptsi.fonctions-deux-variables.b, mpsi-mp2i.fonctions-deux-variables.d, pcsi.fonctions-deux-variables.d, ptsi.fonctions-deux-variables.d, mpsi-mp2i.limites-continuite.c, pcsi.limites-continuite.c, ptsi.limites-continuite.c, bcpst1.derivees-primitives-integrales.b, bcpst1.limites-continuite.c, bcpst1.fonctions-deux-variables.c, mp-mpi.calcul-differentiel-optimisation.f}
\module{Calcul différentiel}
\niveau{L3}
\difficulte{}

\contenu{
\texte{
Soit $f : \Rr 
^2\rightarrow \Rr $ l'application $(x,y) \mapsto 6xy+(y-x)^3.$ On note
$\Delta =\{ (x,y)\in \Rr^2 , -1\leq x\leq y \leq 1\} .$
}
\begin{enumerate}
    \item \question{Dessiner $\Delta $. Montrer que $f$ est born\'ee et atteint ses 
bornes sur $\Delta .$}
    \item \question{Calculer les extrema de $f $ sur le bord de $\Delta $ puis dans 
l'int\'erieur de $\Delta $.}
    \item \question{En d\'eduire les bornes de $f$ sur $\Delta $.}
\end{enumerate}
}
